% Eidesstattliche Erklärung
% Universität der Bundeswehr München

% === Kapitel ohne Nummerierung ===
\chapter*{Eidesstattliche Erklärung}
\phantomsection
\addcontentsline{toc}{chapter}{Eidesstattliche Erklärung}
\thispagestyle{empty}

% === Erklärungstext ===
Hiermit versichere ich eidesstattlich, dass die der Fakultät für Informatik der Universität der Bundeswehr München vorliegende wissenschaftliche Arbeit mit dem Thema:

\begin{center}
  \textit{E-Estonia als Vorbild für die Digitalisierung der deutschen Verwaltung} 
\end{center}

vollkommen selbstständig angefertigt wurde und ich die mit ihr unmittelbar verbundenen Tätigkeiten selbst erbracht habe.

Ich erkläre weiter, dass ich keine anderen als die angegebenen Quellen und Hilfsmittel verwendet habe, insbesondere keine anderen als die angegebenen Informationen. Alle aus gedruckten, ungedruckten oder dem Internet im Wortlaut oder im wesentlichen Inhalt übernommenen Formulierungen und Konzepte sind gemäß den Regeln für wissenschaftliche Arbeiten zitiert und durch genaue Quellenangaben gekennzeichnet.

Die wissenschaftliche Arbeit ist noch nicht anderweitig für Prüfungszwecke und keiner anderen Prüfungsbehörde vorgelegt worden. Sie wurde bisher auch nicht veröffentlicht.

Der Speicherung meiner Arbeit zum Zweck der Plagiatsprüfung stimme ich zu. Ich versichere, dass die elektronische mit der gedruckten Version inhaltlich übereinstimmt.

Hiermit räume ich der Hochschule zum Gebrauch dieser Arbeit in Forschung und Lehre die erforderlichen Nutzungsrechte ein, soweit die Hochschule zur Nutzung durch die gesetzlichen Schrankenbestimmungen des Urhebergesetzes nicht ohne Berechtigung ist. Die Nutzungsrechte verstehen sich als einfache, unentgeltliche, zeitlich, örtlich und inhaltlich unbeschränkte und unwiderrufliche Nutzungsrechte. Die Nutzungsrechte umfassen das Vervielfältigungsrecht, das Verbreitungsrecht, das Ausstellungsrecht und das Recht der öffentlichen Wiedergabe sowie das Recht der Bearbeitung und betreffen alle bekannten und unbekannten Nutzungsarten. Zudem räume ich der Hochschule darüber hinaus das Recht ein, ihrerseits zustimmungsfreies Nutzungsrecht an der Arbeit auf Dritte zu übertragen.

Mir ist bekannt, dass ich verpflichtet bin, die Zustimmung der Hochschule einzuholen, soweit die Verwertung der Arbeit durch mich selbst unter Nennung der Hochschule erfolgen soll.

% === Ort, Datum und Unterschrift ===
\begin{tabbing}
  \hspace*{10cm} \= \hspace*{5cm}\= \kill
  Neubiberg, den \today{} \>
\end{tabbing}

\vspace{-1.5cm}

\begin{tabbing}
  \hspace*{10cm} \= \hspace*{5cm}\= \kill
  \> \textcolor{gray}{..........................................}\\	
  \> \textcolor{gray}{\textit{(Unterschrift)}}
\end{tabbing}